\scanpage{129}
\mspage{5}
% archive foliation: 128
Lorsqu'il n'y a pas de corps de base, il n'est pas clair cependant comment et à quels groupes fondamentaux \deleted{\unsure{respectifs}} pour $X,Y$ on peut rattacher \deleted{par une suite} les groupes fondamentaux des fibres, pour tenir lieu de suite exacte d'homotopie. S'il y a un corps de base, une \deleted{hypothèse} \inserted{conjecture} de départ serait que \deleted{l'image inverse par $f$ de} $\underline\Pi_1(X/Y)$, le pro-schéma en groupes \inserted{$\underline\Pi_1(X)$} sur $X$ qui définit le système local des groupes fondamentaux \deleted{\unsure{en $x$}} de $X$ en ses divers points, et enfin l'image inverse par $f$ du pro-schéma en groupes analogue \inserted{$\underline\Pi_1(Y)$} sur $Y$, sont reliés par une suite exacte \deleted{(qui serait définie d'ailleurs indépendamment de l'existence d'une section de $X$ sur $Y$, car le premier \deleted{\unsure{groupe}} pro-schéma en groupes \inserted{(sur $X$,} nommé serait défini en tous cas). Si $Y$ est le spectre d'un anneau local complet, notamment un anneau de valuation discrète, on devrait pouvoir \unsure{mettre}} il y a de la joie en perspective pour les groupes d'homotopie supérieurs\ldots

Il n'est pas dit que la conjecture envisagée soit plus difficile à démontrer que la partie déjà connue pour les groupes fondamentaux ordinaires. Elle impliquerait que pour des schémas complets $X,Y$ au dessus d'un corps alg clos $k$, on a $\Pi_1(X\times Y/k)=\Pi_1(X/k)\times\Pi_1(Y/k)$. Utilisant tes raisonnements, on en conclut que si $X$ est une variété abélienne sur $k$, alors $\Pi_1(X/k)$ est abélien, et qu'un revêtement principal minimal $X'$ de $X$ est une variété abélienne isogène à $X$. On en déduit
\[
\Pi_1(X/k)=\varprojlim {}_nX
\]
où ${}_nX$ est le noyau (avec sa partie infinitésimale aussi, bien sûr) de la multiplication par $n$, l'homomorphisme de ${}_{mn}X$ dans ${}_nX$ étant induit par la multiplication par $m$. (Utilisant Cartier, ce groupe fondamental est dual, au sens de Cartier, au groupe ind-algébrique limite inductive des ${}_nX^*$, où $X^*$ est la variété duale de $X$). Prenant la $p$-composante de ce groupe fondamental, on trouve ce qui, pour le nombre premier $p$, doit jouer le rôle du module de Weil. Il est hors de doute (et Cartier doit l'avoir fait) qu'il est possible d'associer \deleted{\unsure{à un}} fonctoriellement à un groupe proalgébrique infinitésimal abélien un module sur l'anneau de Witt, et \deleted{\unsure{donc}} qu'il \inserted{\unsure{soit|est}} facile de vérifier qu'en l'occurrence ce module se décompose en les trois morceaux que tu connais bien (correspondants aux trois types principaux de \inserted{$p$-}groupes algébriques abéliens), \deleted{qui} On trouve ainsi
